% !TEX root = ../main.tex
\chapter{Introduction}
\label{ch:intro}

% =====================================================================
% 骨架來源：docs/thesis/thesis-outline.md §1
% 寫作順序：情境 → 現行兩條路都不對 → 我們換一個位置放推論 → 貢獻
% 這一章最後寫。等第 3 章把命題釘死之後再回來。
% =====================================================================

\section{Motivation}
\label{sec:intro:motivation}

% TODO 開場情境：一份模板在實務上會被實例化數十到數千次（不同案場／規格／客戶）。
%      人工填寫耗時且易錯，是自動化的高價值目標。
%      主評測載體是 docx／xlsx 商業表單；DWG 工程圖作為最難的載體獨立呈現。

\section{Two Existing Approaches, Both Wrong}
\label{sec:intro:existing}

% TODO 現行做法一：讓 agent 每次做。問題不只是慢與貴，而是每一次都重新賭一次正確性，
%      而且失敗模式是「沉默失敗」——文件看起來對，某個數字錯了。
% TODO 現行做法二：寫成 skill / SOP 文件。這只是把 prompt 存起來，執行期仍由 LLM 解讀，
%      仍然隨機、仍然吃 token，而且這份 skill 的正確性從來沒有被測試過。
% TODO 【related-work v1 審查要吃下來的三分法，放這裡】
%      RPA = deterministic but manual
%      agents = automatic but stochastic
%      ours = automatic derivation + deterministic execution
%      這是全篇最短的一句定位，不要只放在 related work。

\section{Our Position}
\label{sec:intro:position}

% TODO 把 LLM 用在「一次性推導出規格」，把規格編譯成確定性程式，
%      用可執行 oracle 保證這份規格是對的。O(N)·LLM → O(1)·LLM + O(N)·確定性。
% TODO 【命題已收窄，用這一版，不要用大綱 §0 的舊句】
%      本方法的適用範圍限定於 recurrent template-instantiation tasks whose outputs are
%      deterministically derivable from structured source data and explicit operational rules.
%      收窄的理由見 benchmark-case-v1-thought.md §8.6 與 related-work v1 審查結論。
% TODO 【單一命題，避免連言式新穎性】related-work 審查指出五個 gap 都是
%      「別人各自有、只有我合起來」。重心要收到一句：
%      在什麼可執行條件下，允許把 LLM 從執行路徑上永久移除？

\section{Contributions}
\label{sec:intro:contributions}

% TODO C1–C5，見 thesis-outline.md §0 的貢獻表。
%      注意 C4（zero-LLM execution）被 related-work 審查判定為唯一乾淨的 claim，
%      其他四項都有危險鄰居，論述時要放在 C4 之下而不是並列。

\section{Claims and Where They Are Tested}
\label{sec:intro:claims}

% TODO 三個 claim，每一個都要在第 7 章有對應實驗：
%      1. 正確性不輸逐實例 agent，且執行期變異度為零
%         （但推導期變異度不是零，必須誠實報告，見第 7 章）
%      2. 攤提後成本低一到數個數量級，存在明確 break-even 點
%      3. 失敗會「吵」而不是沉默：模板改版時能被定位、能被修復

\section{Thesis Organization}
\label{sec:intro:organization}

% TODO 最後寫。
